\documentclass[]{article}

\usepackage{amsmath}
\usepackage{multirow}
\usepackage{natbib}
\usepackage{graphicx} 


\begin{document}

\subsection{Simulated Dataset and Experimental Design}
The analysis was performed on a synthetic dataset generated from the numerical integration of the equation of motion for an underdamped harmonic oscillator:
\begin{equation}
m\ddot{x}(t) + b\dot{x}(t) + kx(t) = 0
\end{equation}
where $m$ is the mass (assumed to be unity), $b$ is the damping coefficient, and $k$ is the spring constant. The dataset comprised 20 distinct oscillators, each with unique ground-truth values for $b$ and $k$. For each oscillator, time-series data for position $x(t)$ and velocity $\dot{x}(t)$ were generated over a 20-second interval.

To investigate the limits of parameter observability, the pristine data were systematically degraded. We created an experimental grid by varying two key factors: data fidelity and temporal resolution. Zero-mean Gaussian white noise was injected into both the position and velocity time series to achieve six distinct Signal-to-Noise Ratio (SNR) levels, ranging from 5 dB to 40 dB. Concurrently, the temporal resolution was varied by downsampling the original trajectories to 500, 250, 100, 50, and 25 time steps, while preserving the total 20-second observation window. This process yielded a comprehensive set of test conditions for evaluating the robustness of the estimation algorithms.

\subsection{Parameter Estimation Methodologies}
We compared two distinct methodologies for estimating the parameters $k$ and $b$ from the noisy time-series data, without direct access to acceleration measurements.

\subsubsection{Numerical Derivative with Least-Squares Regression}
This approach approximates the acceleration term $\ddot{x}(t)$ to enable a direct algebraic solution for the parameters. The procedure involved three steps. First, to mitigate the noise amplification inherent in numerical differentiation, the noisy position $x(t)$ and velocity $\dot{x}(t)$ signals were smoothed using a Savitzky-Golay filter. The window size of the filter was dynamically adjusted to be proportional to the temporal resolution of the data. Second, the acceleration $\ddot{x}(t)$ was approximated by applying a central difference formula to the smoothed velocity data. Finally, with the state variables ($x, \dot{x}, \ddot{x}$) estimated, the parameters $b$ and $k$ were determined by solving the rearranged equation of motion using linear least-squares regression:
\begin{equation}
\ddot{x}(t) = - \left(\frac{b}{m}\right) \dot{x}(t) - \left(\frac{k}{m}\right) x(t)
\end{equation}

\subsubsection{Dual Kalman Filter}
As a state-space alternative, we implemented a Dual Kalman Filter (DKF) for simultaneous state and parameter estimation. This method treats the unknown parameters as part of an augmented state vector. The system was modeled with the state vector $\mathbf{z} = [x, \dot{x}, k, b]^T$. The DKF architecture employs two interacting Kalman filters: a state filter estimates the kinematic state ($x, \dot{x}$) based on the current parameter estimates, while a parameter filter estimates the physical parameters ($k, b$) based on the output of the state filter. This iterative, online process allows the model to refine its estimates of both the system's state and its intrinsic parameters at each time step, making it inherently more robust to measurement noise than direct numerical differentiation. The filter was initialized with reasonable prior estimates for the parameters, rather than their ground-truth values, to simulate a realistic application scenario.

\subsection{Performance Evaluation}
To ensure that any observed estimation errors were attributable to noise and temporal sparsity rather than algorithmic flaws, a baseline validation was first conducted. Both methods were applied to the noise-free, high-resolution data, and both successfully recovered the ground-truth parameters with negligible error.

The primary metric for quantifying the accuracy of the parameter estimates was the Mean Absolute Percentage Error (MAPE), calculated for both $k$ and $b$ across all 20 oscillators for each experimental condition. The MAPE is defined as:
\begin{equation}
\text{MAPE} = \frac{1}{N} \sum_{i=1}^{N} \left| \frac{\theta_i - \hat{\theta}_i}{\theta_i} \right| \times 100\%
\end{equation}
where $\theta_i$ is the ground-truth value of a parameter for the $i$-th oscillator, $\hat{\theta}_i$ is its corresponding estimate, and $N$ is the total number of oscillators.

Using this metric, we defined an ``observability threshold'' as the SNR level below which the MAPE for a given parameter exceeded 5\%. This threshold serves as a quantitative boundary, delineating the minimum data fidelity required for reliable parameter recovery with each method.

\end{document}
                